\begin{table}[h!]
\caption{HGDTI~\cite{yu2022hgdti} rerun with its own released pipeline (model module rebuilt from the released bytecode), AUROC as the mean over three seeds of the ten-fold mean (seed SD $\le 0.008$). Each step adds one correction to the previous one; steps 3--4 score the step-2 model. Floor: best model-free score on the same test pairs (\texttt{protocheck} 0.2.0; in step 1 the propagated adjacency holds every test edge, so the floor is the input handed back; steps 2--3 protein degree, step 4 Adamic--Adar). Random-trained: the same folds, positives and test sets, with training and validation negatives drawn uniformly from unknown pairs instead of the paper's similarity-filtered ``reliable'' ones. The released pipeline run with uniform negatives, the ``random sampling'' bar of the paper's Fig.~8 (0.953), gives 0.9976.}
\label{tab:hgdti_rerun}
\footnotesize
\resizebox{\textwidth}{!}{%
\begin{tabular}{lccc}
\hline
 & \multicolumn{2}{c}{Reliable negatives (as released)} & Random-trained \\
\cline{2-3}
Step & HGDTI & Floor & HGDTI \\
\hline
1. As released (adjacency from the test labels, reliable test negatives) & 0.9995 & 1.000 & -- \\
2. Propagation over training edges only & 0.9462 & 0.8865 & 0.9234 \\
3. Uniform test negatives & 0.8801 & 0.8845 & 0.9016 \\
4. Degree-matched test negatives & 0.5728 & 0.7395 & 0.6631 \\
\hline
\end{tabular}
}
\end{table}
